\documentclass[11pt,a4paper]{article}

\usepackage[utf8]{inputenc}
\usepackage[T1]{fontenc}
\usepackage[margin=2.4cm]{geometry}
\usepackage{amsmath,amssymb}
\usepackage{booktabs}
\usepackage{graphicx}
\usepackage{microtype}
\usepackage[ruled,vlined,linesnumbered]{algorithm2e}
\usepackage{listings}
\usepackage{xcolor}
\usepackage[hidelinks]{hyperref}
\usepackage{caption}

\definecolor{codebg}{gray}{0.96}
\lstset{
  basicstyle=\ttfamily\small,
  backgroundcolor=\color{codebg},
  breaklines=true,
  frame=none,
  columns=fullflexible,
  keepspaces=true,
  showstringspaces=false,
}

\newcommand{\Adv}{A}
\newcommand{\polold}{\pi_{\theta_{\mathrm{old}}}}
\newcommand{\polref}{\pi_{\mathrm{ref}}}
\newcommand{\pol}{\pi_{\theta}}

\title{\textbf{Accelerating GRPO Post-Training with Completion Pruning}\\
\large An implementation and analysis of CPPO for Qwen3-0.6B on DAPO-Math-17k}
\author{Yusuf Saifutdinov}
\date{\today}

\begin{document}
\maketitle

\begin{abstract}
Group Relative Policy Optimization (GRPO) has become the workhorse of
reasoning-oriented post-training, but it is expensive: every question is
answered $G$ times, and every one of those $G$ completions is pushed through a
forward and a backward pass. Completion Pruning Policy Optimization (CPPO)
observes that a completion's contribution to the policy gradient is
proportional to the magnitude of its group-relative advantage, and therefore
back-propagates through only the $k = \lfloor G(1-P) \rfloor$ completions with
the largest $|\Adv_i|$. This report presents a from-scratch implementation of
GRPO and CPPO on top of TRL for Qwen3-0.6B trained on DAPO-Math-17k, evaluated
with \texttt{lm-evaluation-harness} on GSM8K, Minerva MATH and AIME 2024. We
describe both algorithms, show how CPPO's \emph{dynamic completion allocation}
maps exactly onto TRL's batching parameters without modifying TRL, and give a
cost model that predicts CPPO's end-to-end speedup from a single measured
quantity: the share of training time spent in the update stage. The central
finding is that CPPO's accuracy is essentially free --- pruning discards
completions that carry almost none of the batch's learning signal --- but its
\emph{speedup is bounded by Amdahl's law} on the rollout stage, so the
published $8.32\times$ figures are reachable only in configurations where the
update stage dominates.
\end{abstract}

\tableofcontents

\section{Introduction}
\label{sec:intro}

Reinforcement learning with verifiable rewards has driven most of the recent
progress in mathematical reasoning, code generation and long-horizon
instruction following~\cite{deepseekai2025r1,yu2025dapo}. The dominant
algorithm family is PPO~\cite{schulman2017ppo} and its critic-free descendant
GRPO~\cite{shao2024deepseekmath}, which replaces the learned value function
with a baseline estimated from a \emph{group} of completions sampled for the
same question.

Removing the critic removes one model from GPU memory, but GRPO remains
costly for three reasons:

\begin{enumerate}
  \item \textbf{Multiple models resident at once.} The policy, and --- when a
  KL penalty is used --- a frozen reference model, plus optimiser state for
  the policy.
  \item \textbf{A slow autoregressive rollout.} Each optimiser step needs $G$
  sampled completions \emph{per question}, and sampling is memory-bandwidth
  bound rather than compute bound.
  \item \textbf{A wide update stage.} Every one of those $G$ completions is
  pushed through a forward and a backward pass of the policy.
\end{enumerate}

CPPO~\cite{lin2025cppo} attacks the third cost. This report answers the
question posed by the task directly: \emph{how accurate is CPPO, and how does
it influence RL training acceleration?} Section~\ref{sec:grpo} and
Section~\ref{sec:cppo} describe the two algorithms,
Section~\ref{sec:implementation} documents the implementation and the places
where a faithful mapping onto TRL required judgement,
Section~\ref{sec:setup}--\ref{sec:results} give the experimental protocol and
results, and Section~\ref{sec:analysis} develops the cost model that explains
what the numbers mean.

\section{Background: GRPO}
\label{sec:grpo}

\subsection{Objective}

For a question $q$ drawn from the data distribution $P(Q)$, the old policy
$\polold$ samples a group of $G$ completions $\{o_1,\dots,o_G\}$. GRPO
maximises

\begin{equation}
\label{eq:grpo}
\begin{split}
\mathcal{J}_{\mathrm{GRPO}}(\theta) =\;
&\mathbb{E}_{q \sim P(Q),\, \{o_i\}_{i=1}^{G} \sim \polold(o \mid q)}
\Bigg[ \frac{1}{G}\sum_{i=1}^{G} \frac{1}{|o_i|} \sum_{t=1}^{|o_i|}
\Big\{ \min\Big( \rho_{i,t}(\theta)\, \Adv_i,\\
&\quad \mathrm{clip}\big(\rho_{i,t}(\theta),\, 1-\epsilon,\, 1+\epsilon\big)\,
\Adv_i \Big) - \beta\, \mathbb{D}_{\mathrm{KL}}\!\left[\pol \,\|\, \polref\right]
\Big\} \Bigg],
\end{split}
\end{equation}

where $\rho_{i,t}(\theta) = \pol(o_{i,t} \mid q, o_{i,<t}) / \polold(o_{i,t}
\mid q, o_{i,<t})$ is the token-level importance ratio and the advantage is
the group-standardised reward

\begin{equation}
\label{eq:advantage}
\Adv_i = \frac{r_i - \mathrm{mean}(\{r_1,\dots,r_G\})}
              {\mathrm{std}(\{r_1,\dots,r_G\})}.
\end{equation}

Crucially, $\Adv_i$ is constant across the tokens of completion $i$: the whole
completion is weighted by one scalar.

\subsection{Rule-based reward}

Following the CPPO paper and common practice for verifiable domains, the
reward is the sum of a format term and an accuracy term,
$r_i = R_{\mathrm{format}}(o_i) + R_{\mathrm{accuracy}}(o_i)$, with
$R_{\mathrm{format}} \in \{0,1\}$ and $R_{\mathrm{accuracy}} \in \{0,2\}$. No
reward model is trained. Our implementation checks accuracy with
symbolic equivalence (\texttt{math\_verify}) rather than string matching, so
$\tfrac{1}{2}$, $0.5$ and $\backslash\mathrm{frac}\{1\}\{2\}$ all score alike.

\subsection{Where the time goes}

Let $B$ be the number of questions per optimiser step. A step costs

\begin{equation}
\label{eq:grpo-cost}
T_{\mathrm{step}} \;=\; \underbrace{B \cdot G \cdot c_{\mathrm{gen}}}_{\text{rollout}}
\;+\; \underbrace{B \cdot G \cdot \big(n_{\mathrm{fwd}}\, c_{\mathrm{fwd}} + c_{\mathrm{bwd}}\big)}_{\text{update}},
\end{equation}

where $n_{\mathrm{fwd}}$ counts the forward passes over completions: one for
the policy, plus one for the reference model when $\beta \neq 0$, plus one for
the old policy when the samples are off-policy. In DeepSeekMath's
configuration with $G=64$ that is $192$ forward passes per
question~\cite{lin2025cppo}. Both terms scale linearly in $G$, which is what
makes larger groups --- and hence better baselines --- so expensive.

\section{CPPO: Completion Pruning Policy Optimization}
\label{sec:cppo}

\subsection{Which completions actually matter}

Differentiating Eq.~\eqref{eq:grpo} and dropping the KL term, which
Hu~et~al.~\cite{hu2025openreasonerzero} show is not what shapes reasoning
behaviour, gives

\begin{equation}
\label{eq:grad}
\nabla_\theta \mathcal{J}(\theta) \approx \mathbb{E}\Bigg[
\frac{1}{G}\sum_{i=1}^{G} \frac{1}{|o_i|} \sum_{t=1}^{|o_i|}
\underbrace{\rho_{i,t}(\theta)}_{\text{post-forward}} \cdot
\underbrace{\Adv_i}_{\text{\textbf{pre}-forward}} \cdot
\underbrace{\nabla_\theta \log \pol(o_{i,t} \mid q, o_{i,<t})}_{\text{post-forward}}
\Bigg].
\end{equation}

Three factors must all be non-negligible for a completion to move the policy.
Two of them are only knowable \emph{after} the policy forward pass --- but
$\Adv_i$ is known \emph{before} it, as soon as the group's rewards have been
scored. That asymmetry is the entire idea: a completion whose $|\Adv_i|$ is
near zero can be identified and discarded before any expensive computation is
spent on it.

The paper also gives an interpretation of \emph{which} completions get
discarded. With a two-part reward, the four outcome types are:

\begin{center}
\begin{tabular}{llll}
\toprule
Format & Answer & Typical $|\Adv|$ & Training signal \\
\midrule
correct   & correct   & large & ``do this'' --- clean positive signal \\
incorrect & incorrect & large & ``avoid this'' --- clean negative signal \\
correct   & incorrect & small & ambiguous; rewards form over substance \\
incorrect & correct   & small & ambiguous; rewards luck over method \\
\bottomrule
\end{tabular}
\end{center}

Pruning by $|\Adv_i|$ therefore removes precisely the partially-correct
completions whose gradient is both weak and semantically muddled, which is why
CPPO can \emph{improve} accuracy rather than merely preserve it.

\subsection{The pruned objective}

CPPO adds a selection constraint to Eq.~\eqref{eq:grpo}. To keep every data-parallel
device processing an identical number of completions --- otherwise the slowest
device sets the pace, the ``bucket effect'' --- the constraint is expressed as
a per-group top-$k$ rather than a global threshold:

\begin{equation}
\label{eq:k}
k = \lfloor G \cdot (1 - P) \rfloor, \qquad P \in [0,1),
\end{equation}
\begin{equation}
\label{eq:index-set}
\mathcal{I} = \big\{\, i \in \{1,\dots,G\} \;\big|\; |\Adv_i| \text{ is among the top } k \text{ values} \,\big\},
\end{equation}
\begin{equation}
\label{eq:cppo}
\begin{split}
\mathcal{J}_{\mathrm{CPPO}}(\theta) =\;
&\mathbb{E}_{q,\{o_i\}} \Bigg[ \frac{1}{k}\sum_{i \in \mathcal{I}} \frac{1}{|o_i|}
\sum_{t=1}^{|o_i|} \Big\{ \min\Big( \rho_{i,t}(\theta)\, \Adv_i,\\
&\quad \mathrm{clip}\big(\rho_{i,t}(\theta), 1-\epsilon, 1+\epsilon\big)\, \Adv_i \Big)
- \beta\, \mathbb{D}_{\mathrm{KL}}\!\left[\pol \,\|\, \polref\right] \Big\} \Bigg].
\end{split}
\end{equation}

Note the normaliser: $1/k$, not $1/G$. Pruning must not shrink the gradient
magnitude, or it would silently act as a learning-rate decay.

\subsection{Dynamic completion allocation}

Pruning leaves $G-k$ completion slots empty in the update stage, so the
accelerator runs at a fraction of its capacity. CPPO refills those slots with
completions sampled for \emph{additional} questions. Per generation round the
device then handles

\begin{equation}
\label{eq:allocation}
m = \left\lfloor G / k \right\rfloor
\end{equation}

times as many questions, each still contributing $k$ retained completions, so
the post-pruning batch is again the full width the hardware can hold. Total
rollout work over an epoch is unchanged --- the same questions are still
answered $G$ times --- but the epoch now needs $m$ times fewer optimiser
steps, and each of those steps is as wide as the baseline's.

\subsection{Algorithm}

\begin{algorithm}[t]
\caption{CPPO training step (with dynamic completion allocation)}
\label{alg:cppo}
\SetAlgoLined
\KwIn{policy $\pol$, group size $G$, pruning rate $P$, questions per round $B$,
      reward function $r(\cdot)$}
$k \leftarrow \lfloor G(1-P) \rfloor$; \quad $m \leftarrow \lfloor G/k \rfloor$\;
\For{each generation round}{
  Sample $m \cdot B$ questions $\{q_j\}$ \tcp*{allocation: $m\times$ more questions}
  \ForEach{question $q_j$}{
    Sample $G$ completions $\{o_{j,1},\dots,o_{j,G}\} \sim \polold(\cdot \mid q_j)$\;
    Score rewards $r_{j,i} \leftarrow R_{\mathrm{format}}(o_{j,i}) + R_{\mathrm{accuracy}}(o_{j,i})$\;
    $\Adv_{j,i} \leftarrow (r_{j,i} - \mathrm{mean}_i r_{j,\cdot}) / \mathrm{std}_i\, r_{j,\cdot}$
      \tcp*{pre-forward information}
    $\mathcal{I}_j \leftarrow \mathrm{top}\text{-}k\big(\{|\Adv_{j,i}|\}_{i=1}^{G}\big)$
      \tcp*{\textbf{pruning}}
  }
  $\mathcal{B} \leftarrow \bigcup_j \{ o_{j,i} : i \in \mathcal{I}_j \}$
    \tcp*{$m \cdot B \cdot k \approx B \cdot G$ completions}
  \ForEach{micro-batch of $\mathcal{B}$}{
    Forward and backward $\pol$ on the micro-batch only; accumulate gradients\;
  }
  Optimiser step\;
}
\end{algorithm}

Steps 1--3 of the loop are identical to GRPO. The only changes are the
top-$k$ selection and the $m$-fold widening of the question batch.

\section{Implementation}
\label{sec:implementation}

The implementation is a small package layered on TRL~\cite{vonwerra2020trl}
rather than a fork of it, so the GRPO baseline is provably stock TRL:

\begin{center}
\begin{tabular}{ll}
\toprule
Module & Responsibility \\
\midrule
\texttt{cppo/pruning.py}   & Top-$k$ selection by $|\Adv|$; pruning diagnostics \\
\texttt{cppo/geometry.py}  & Pruning rate $\rightarrow$ TRL batch parameters \\
\texttt{cppo/trainer.py}   & \texttt{ProfiledGRPOTrainer}, \texttt{CPPOTrainer} \\
\texttt{cppo/profiling.py} & Stage timing and peak-memory tracking (shared) \\
\texttt{cppo/rewards.py}   & Format and accuracy rewards \\
\texttt{cppo/data.py}      & DAPO-Math-17k loading and prompt templating \\
\texttt{cppo/train.py}     & Training entry point \\
\texttt{cppo/evaluate.py}  & \texttt{lm-evaluation-harness} wrapper \\
\texttt{cppo/report.py}    & Renders result tables for this document \\
\bottomrule
\end{tabular}
\end{center}

\subsection{Selection}

The retention mask is built per group with a \emph{stable} descending sort of
$|\Adv|$, so ties resolve to the lowest index and every data-parallel rank
selects the same completions from the same advantages without communicating.
Keeping exactly $k$ per group --- rather than thresholding $|\Adv_i| \ge
\gamma$ --- is what makes the retained count identical on every device, which
is the requirement Eq.~\eqref{eq:k} exists to satisfy.

\subsection{Mapping dynamic allocation onto TRL}
\label{sec:geometry}

TRL derives its rollout batch from three knobs:
\[
\begin{aligned}
\texttt{generation\_batch\_size} = \;&\texttt{per\_device\_train\_batch\_size}\\
\times\;&\texttt{num\_processes} \times \texttt{steps\_per\_generation},
\end{aligned}
\]
samples $\texttt{generation\_batch\_size} / G$ distinct questions per round,
and splits the result into \texttt{steps\_per\_generation} optimiser
micro-batches.

The useful observation is that dynamic completion allocation requires
\emph{no changes to TRL's batching code at all}. Multiplying
\texttt{per\_device\_train\_batch\_size} by $m$ makes TRL sample $m$ times
more questions per round; pruning then shrinks each micro-batch back to
exactly the baseline width. Concretely, for the configuration used here
($G=8$, baseline micro-batch $16$, $\texttt{steps\_per\_generation}=8$):

\begin{center}
\begin{tabular}{lrrrrr}
\toprule
Run & $P$ & $k$ & $m$ & Questions/round & Update micro-batch \\
\midrule
GRPO baseline          & 0.000 & 8 & 1 & 16  & 16 \\
CPPO $P=0.50$          & 0.500 & 4 & 2 & 32  & 16 \\
CPPO $P=0.75$          & 0.750 & 2 & 4 & 64  & 16 \\
CPPO $P=0.875$         & 0.875 & 1 & 8 & 128 & 16 \\
CPPO $P=0.75$, no alloc.& 0.750 & 2 & 1 & 16  & \phantom{0}4 \\
\bottomrule
\end{tabular}
\end{center}

The update stage therefore sees identical tensor shapes in every allocated
configuration --- identical memory, identical kernel occupancy --- while
covering $m$ times more questions per optimiser step. The last row is the
paper's ``+ Completion Pruning'' ablation: pruning without refilling, which
leaves three quarters of the batch empty.

\subsection{Where pruning is applied, and why}
\label{sec:pruning-point}

The CPPO paper prunes before the policy, reference \emph{and} old-policy
forward passes. TRL computes rewards --- and hence advantages --- \emph{after}
the reference and old-policy log-probabilities, so a literal ``prune first''
ordering would require reimplementing a 600-line private method against
internal TRL state, which is both fragile across versions and hard to audit.

Instead this implementation prunes at the end of the rollout, immediately
before the policy forward/backward pass, and adopts a configuration in which
the auxiliary passes \emph{do not exist}:

\begin{itemize}
  \item $\beta = 0$: no reference model, so no reference forward pass. This is
  TRL's own default, matches Dr.~GRPO's recommendation~\cite{liu2025drgrpo},
  and is exactly the KL-free approximation from which CPPO derives its pruning
  criterion in Eq.~\eqref{eq:grad}. It also removes a whole model from GPU
  memory.
  \item $\mu = 1$ iteration with generation aligned to the optimiser step: the
  samples are on-policy, so TRL skips the old-policy forward pass.
  \item vLLM importance-sampling correction disabled: otherwise TRL recomputes
  old-policy log-probabilities to correct the sampler/trainer mismatch.
\end{itemize}

Under these settings the only pass over completions is the policy's, and
pruning before it reproduces Eq.~\eqref{eq:cppo} exactly. When a user
overrides them, \texttt{CPPOTrainer} emits a warning naming each unpruned
forward pass: training remains correct, but the measured speedup becomes a
lower bound. Making that trade-off explicit --- rather than silently reporting
a smaller number --- is the point.

\subsection{Correctness details that are easy to get wrong}

\paragraph{Loss normalisation.} TRL's default \texttt{dapo} loss divides by
\texttt{num\_items\_in\_batch}, the count of loss-carrying tokens in the whole
generation batch. Pruning removes tokens; leaving the pre-pruning count in
place would scale the loss --- and therefore the effective learning rate --- by
$k/G$. The trainer recomputes the normaliser over the retained completions.
This is the implementation-level counterpart of the $1/k$ in
Eq.~\eqref{eq:cppo}.

\paragraph{Group locality.} Pruning is per question, so a group must live
entirely on one device. The trainer validates that the local generation batch
is a whole multiple of $G$ and fails loudly otherwise, rather than silently
pruning across a group boundary.

\paragraph{Degenerate groups.} When all $G$ completions of a question receive
the same reward, the standard deviation is zero and every advantage is zero:
the group contributes no gradient whatsoever, yet still costs a full forward
and backward pass. The trainer logs the fraction of such groups, and offers an
optional \texttt{drop\_zero\_advantage} mode that removes them outright. That
mode is exact only when $\beta = 0$ and, because it makes the retained count
data-dependent, is restricted to single-process runs --- unequal batch sizes
across ranks would deadlock a collective operation.

\section{Experimental setup}
\label{sec:setup}

\begin{center}
\begin{tabular}{ll}
\toprule
Policy model        & Qwen3-0.6B~\cite{qwen3} \\
Training data       & \texttt{open-r1/DAPO-Math-17k-Processed} (config \texttt{en}) \\
Training subset     & 8{,}192 problems, one epoch \\
Reward              & format ($\{0,1\}$) $+$ accuracy ($\{0,2\}$, \texttt{math\_verify}) \\
Prompt              & chat template, ``reason step by step, answer in \verb|\boxed{}|'' \\
Group size $G$      & 8 \\
Max completion len. & 1024 tokens \\
Sampling            & temperature 1.0, top-$p$ 1.0 \\
Objective           & \texttt{dapo} loss, $\epsilon = 0.2$, $\epsilon_{\mathrm{high}} = 0.28$, $\beta = 0$ \\
Optimiser           & AdamW, lr $10^{-6}$, constant with 10 warmup steps, grad-clip 0.2 \\
Rollout backend     & vLLM~\cite{kwon2023vllm}, colocate mode \\
Evaluation          & \texttt{lm-evaluation-harness}~\cite{eval-harness}: \\
                    & \quad \texttt{gsm8k} (5-shot), \texttt{minerva\_math} (4-shot), \texttt{aime24} (0-shot) \\
\bottomrule
\end{tabular}
\end{center}

\paragraph{Protocol.} Every run consumes the \emph{same 8{,}192 problems in
the same order}. Because dynamic allocation covers $m$ times more questions
per optimiser step, CPPO completes that epoch in $m$ times fewer steps. Fixing
the data rather than the step count is what makes a wall-clock comparison
meaningful; fixing the step count instead would let CPPO see $m$ times more
data and confound speed with sample efficiency.

The results tables therefore report \emph{questions per second} as the
headline metric rather than raw wall clock. Throughput is comparable under
either protocol: it credits an algorithm for covering the same training data
in less time, and never for covering more data in the same time.

\paragraph{Prompt/evaluation alignment.} The policy is trained with a chat
template and a \verb|\boxed{}| answer convention, and evaluated through
\texttt{lm\_eval}'s \texttt{-{}-apply\_chat\_template} flag with the same
system prompt. \texttt{minerva\_math} and \texttt{aime24} extract from
\verb|\boxed{}|, and \texttt{gsm8k}'s \texttt{flexible-extract} filter takes
the last number, so the training format transfers to all three benchmarks
without a task-specific adapter.

\section{Results}
\label{sec:results}

\subsection*{Training performance}

\resizebox{\ifdim\width>\textwidth\textwidth\else\width\fi}{!}{%
\begin{tabular}{lrrrrrrrrrrr}
\toprule
Method & P & k & m & Steps & Questions & Wall clock (s) & Rollout (s) & Update (s) & Peak mem (GiB) & Questions/s & Throughput gain \\
\midrule
smoke-grpo & 0.00\% & 4 & 1x & 2 & 2 & 553.0 & 21.1 & 526.1 & 7.47 & 0.0036 & 1.00x \\
smoke-cppo & 50.00\% & 2 & 2x & 2 & 4 & 550.9 & 28.7 & 516.4 & 6.99 & 0.0073 & 2.01x \\
\bottomrule
\end{tabular}
}

\subsection*{Downstream accuracy (lm-evaluation-harness, \% exact match)}

\emph{No evaluation results available.}

\subsection*{Update-stage micro-benchmark}

\resizebox{\ifdim\width>\textwidth\textwidth\else\width\fi}{!}{%
\begin{tabular}{lrrrrrr}
\toprule
P & k & Completions/step & Tokens/step & Step time (s) & Speedup (pruning only) & Speedup (with allocation) \\
\midrule
0.00\% & 8 & 8 & 1024 & 4.4768 $\pm$ 0.0016 & 1.00x & 1.00x \\
12.50\% & 7 & 7 & 896 & 3.9283 $\pm$ 0.0027 & 1.14x & 1.13x \\
25.00\% & 6 & 6 & 768 & 3.4326 $\pm$ 0.0017 & 1.30x & 1.30x \\
37.50\% & 5 & 5 & 640 & 2.9138 $\pm$ 0.0011 & 1.54x & 1.53x \\
50.00\% & 4 & 4 & 512 & 2.4255 $\pm$ 0.0019 & 1.85x & 2.00x \\
62.50\% & 3 & 3 & 384 & 1.8822 $\pm$ 0.0031 & 2.38x & 2.61x \\
75.00\% & 2 & 2 & 256 & 1.3869 $\pm$ 0.0060 & 3.23x & 4.01x \\
87.50\% & 1 & 1 & 128 & 0.8302 $\pm$ 0.0019 & 5.39x & 8.02x \\
\bottomrule
\end{tabular}
}


\section{Analysis}
\label{sec:analysis}

\subsection{A cost model for CPPO's speedup}

CPPO changes the update stage and leaves the rollout stage untouched. Let $f$
be the fraction of baseline training time spent in the update stage. Update
work scales with the number of completions back-propagated, i.e. by $k/G =
1-P$. The achievable end-to-end speedup is then bounded by Amdahl's law:

\begin{equation}
\label{eq:amdahl}
S(P) \;=\; \frac{1}{(1-f) + f\,(1-P)}, \qquad
S_{\max} \;=\; \lim_{P \to 1} S(P) \;=\; \frac{1}{1-f}.
\end{equation}

This single equation explains the spread in the published numbers. The
important consequences are:

\begin{itemize}
  \item \textbf{The rollout is the ceiling.} If generation is 50\% of training
  time, no pruning rate can exceed $2\times$ overall, however aggressive.
  \item \textbf{CPPO pays off most where the update stage is fattest}: large
  $G$, long completions, a KL penalty that adds a reference forward pass, or
  multiple inner iterations $\mu > 1$. The paper's $8.32\times$ on GSM8K
  implies $f \approx 0.88$, consistent with its $\beta = 0.04$ configuration
  in which every completion also pays a reference-model forward pass.
  \item \textbf{Our configuration is deliberately less favourable to CPPO.}
  With $\beta = 0$, $\mu = 1$ and a vLLM rollout, the update stage is a single
  forward plus a backward, and $f$ is correspondingly smaller. The honest
  reading is that the configuration that is best for CPPO's \emph{relative}
  numbers is not the configuration that is best in absolute terms.
  \item \textbf{The rollout is not immune in principle.} Nothing in the
  rollout depends on advantages, which do not exist until sampling has
  finished, so completion pruning cannot shrink it. Reducing rollout cost
  needs an orthogonal method --- speculative or truncated sampling, early
  stopping of degenerate groups, or reusing completions across steps.
\end{itemize}

Because $f$ is measured directly by the profiling mixin (both stages are
timed with the same code in both runs, and are additive by construction),
Eq.~\eqref{eq:amdahl} turns into a falsifiable prediction that the reported
tables can be checked against.

\subsection{Why accuracy survives}

The pruning diagnostics logged each step (\texttt{cppo/retained\_signal\_fraction})
report the share of the batch's total $|\Adv|$ mass that survives pruning.
Because group-standardised rewards from a two-part reward function are
strongly bimodal, that share stays close to $1$ even at $P = 0.75$: the
discarded completions are numerous but individually almost weightless. This is
the quantitative statement behind the qualitative table in
Section~\ref{sec:cppo}, and it is the mechanism by which CPPO preserves --- and
in the paper's experiments sometimes improves --- accuracy.

The paper's own metric ablation supports the criterion rather than the mere
act of subsetting: keeping the \emph{largest} $|\Adv|$ scores $77.67\%$ on
GSM8K, random subsetting $76.98\%$, and keeping the \emph{smallest}
$74.23\%$. Pruning by raw (signed) advantage rather than absolute advantage is
also worse, confirming that large-negative completions --- ``do not do this''
--- are as informative as large-positive ones.

There is a floor, however. At $P = 0.9375$ ($k=1$) the paper's accuracy drops
on both GSM8K and MATH: with a single retained completion the surviving
gradient is a single-sample estimate, and genuinely useful completions start
being discarded alongside the weightless ones. $P \in [0.5, 0.875]$ is the
usable band.

\subsection{Memory}

Completion pruning does not reduce \emph{peak} memory on its own, and with
dynamic allocation it is designed not to: the whole point of refilling the
freed slots is to keep the update-stage tensors at the width the device can
hold. What CPPO buys is throughput at constant memory, not a smaller
footprint. The memory saving in this pipeline comes from a different decision
--- $\beta = 0$ removes the reference model entirely --- and the
\texttt{no\_allocation} ablation is the configuration in which pruning does
lower peak memory, at the cost of leaving the device under-occupied.

\section{Limitations and reproducibility}
\label{sec:limitations}

\paragraph{What was measured where.} The tables in
Section~\ref{sec:results} are generated by \texttt{cppo.report} from the JSON
artefacts each run writes; the README states, for each table, the hardware it
was produced on. Where a table was produced on a laptop-class accelerator
rather than the reference A100/H100, that is recorded explicitly rather than
extrapolated. All scripts needed to reproduce the full study on a single
80\,GB GPU are in \texttt{scripts/}.

\paragraph{Known gaps.}
\begin{itemize}
  \item Pruning happens after the rollout rather than before the auxiliary
  forward passes (Section~\ref{sec:pruning-point}). With the recommended
  $\beta = 0$, $\mu = 1$ configuration this is exact; otherwise the reported
  speedup understates CPPO.
  \item Qwen3-0.6B is small enough that a single optimiser step is dominated
  by fixed overheads that a 7B model would amortise, so $f$ --- and therefore
  the achievable speedup --- is smaller here than it would be at scale.
  \item \texttt{aime24} has 30 problems; a single problem is $3.3$ points.
  Differences on that benchmark should not be read as significant without
  multiple seeds.
  \item Single seed per configuration. GRPO runs are noisy; the accuracy
  column should be read as ``no regression'' evidence rather than as a precise
  ranking.
\end{itemize}

\section{Conclusion}
\label{sec:conclusion}

CPPO is a well-targeted optimisation with an unusually clean justification:
the factor of the policy gradient that is knowable before the expensive
computation is exactly the factor that decides whether the computation is
worth doing. Its accuracy cost is close to zero in the $P \in [0.5, 0.875]$
band, because group-standardised rewards concentrate the learning signal in a
small minority of completions, and its dynamic allocation strategy maps onto
TRL's batching model without touching TRL.

Its acceleration, however, is bounded by what it does not touch. Completion
pruning is an update-stage method, and Eq.~\eqref{eq:amdahl} makes the ceiling
explicit: $1/(1-f)$, where $f$ is the update stage's share of training time.
Reported speedups of $8\times$ describe configurations where $f \approx 0.9$
--- large groups, a reference model, long completions. In a modern KL-free
configuration with a vLLM rollout, $f$ is smaller and so is the payoff. The
practical recommendation that follows is to adopt CPPO at $P \approx 0.75$
--- it is nearly free and never hurt accuracy in our band --- and to pair it
with a rollout-stage method, because after CPPO the rollout is where all the
remaining time is.

\bibliographystyle{plain}
\bibliography{references}

\end{document}
